\documentclass[11pt]{article}
\usepackage[margin=2.5cm]{geometry}
\usepackage{enumitem,booktabs,hyperref}
\setlist{itemsep=2pt,topsep=3pt}
\newcommand{\rev}[1]{\medskip\noindent\textbf{#1}}
\newcommand{\q}[1]{\par\smallskip\noindent\emph{#1}\par\smallskip}
\title{Response to the reviews of CAOR-D-26-02865 and summary of changes\\
\large Manuscript prepared for the \emph{European Journal of Operational Research}}
\author{Quang-Vinh Dang}
\date{3 November 2026}
\begin{document}
\maketitle

This note lists every point raised in the four reviews, what was changed in the
manuscript (\texttt{paper/cascade-mis.tex}), where, and what was \emph{not}
done. Section and table numbers refer to the new manuscript. The paper was
rejected; the manuscript is rewritten for a new submission, and the note is
written so that each change can be checked against the text and the data in
\texttt{results/}.

\section*{Summary}
\begin{itemize}
\item \textbf{The central construction is not new} (Reviewers 1, 3). The region
  $R(F)$ is the construction of Borowitz, Gro{\ss}mann and Schulz (ALENEX 2025). The
  paper no longer claims it. It is now presented as the use of that construction
  as the inner move of a \emph{static} solver, with reductions applied inside the
  region and an online rule for $|F|$, and the contribution list says so.
\item \textbf{The result tables were wrong} (Reviewers 1, 4). The quality,
  aggregate and budget tables had been generated from five of the seven solvers,
  and the Wilcoxon test had been run on raw sizes although the text said relative
  gaps. Both are fixed at the source scripts. The corrected statistics change the
  claims: the method is significantly ahead of three of six baselines, not five,
  and third of seven on mean relative gap.
\item \textbf{New experiments} (own solver only): residual-kernel statistics and
  whole-input exact solves; exact against heuristic repair and regional
  reductions on/off; a sweep over the initial release size; independent Delaunay
  graphs.
\item \textbf{Not done}: no recent baseline (InfVC) was run, and no baseline was run
  on the new Delaunay graphs. See ``What remains open''.
\end{itemize}

\section*{Reviewer 3}

\rev{3.1 Explain the Twin, Unconfined and LP rules.}
Section 2 now has a paragraph for each, with the condition, the effect, why it is
safe (a two-line exchange argument for the twin rule with an edge in the shared
neighbourhood), and the implementation bounds on the unconfined test (set size 8,
$2\times10^4$ scanned edges; a capped test answers ``confined'', so the cap
cannot make the rule unsound).

\rev{3.2 How does $R(F)$ differ from the region of Borowitz et al.?}
It does not, in essence. After closing their subgraph $H$ under incumbent
neighbours and adding the tight vertices, their subgraph is $R(F)$ for $F$ the
incumbent vertices it contains. A Remark after Theorem 5 (the region theorem) says so, the
introduction states it, and Table 1 compares the two constructions along six
axes (problem, how the released set is chosen, candidate restriction,
sub-problem objective, reductions, sub-solver, size control). What differs is
the setting (static against dynamic), the released set (count of incumbent
vertices along a walk, against BFS depth on all vertices), no high-degree pruning
here, an in-house reusable sub-solver, and an online rule for $|F|$ in place of
fixed $d,\nu_{\max}$. I did not run their code, so the comparison is of
constructions, not of performance, and the paper says that.

\rev{3.3 Algorithm 1 is unclear.}
Algorithm 1 is rewritten with explicit input and output, and a new parameter table (Section 4.1)
gives every parameter and its value: (1) initial $k_0=4$; doubling after
$s=512$ consecutive proved sub-solves without gain; halving when a sub-solve is
cut off or the region exceeds 3000 vertices; $+\lfloor k/4\rfloor+1$ after a
sweep with no gain. (2) $F$ is the first $k$ incumbent vertices met by a
breadth-first walk from the seed, visiting at most $\max(64,64k)$ vertices.
(3) Plateau moves switch on after one full sweep without gain, and off after a
sweep with gain. (4) A perturbation runs after every batch of 64 region moves;
a restart fires after 200 failed full-strength kicks or after 200 first-dive
durations without improvement. (5) Input: graph, budget, seed; output: a
maximal independent set. All values are read from the code, not from memory of
the design.

\rev{3.4, 3.5 More recent baselines; update related work.}
Related work now covers Borowitz et al.\ (2025), Hashimoto et al.\ (2019),
Katzmann and Komusiewicz (2017), InfVC (IJCAI 2025), and the 2024 weighted-case
survey and concurrent local search of Gro{\ss}mann, Langedal and Schulz. \textbf{No
recent solver was run}; Section 5.15 says that the ranking claims are claims
against 2012--2019 solvers.

\rev{3.6 Sensitivity to $k$.}
New table (Section 5.9): the initial value $k_0\in\{1,4,16,64,256,1024\}$ on
four instances. On \texttt{del18} the results span 0.023\% (the seed spread is
0.024\%), and the final $k$ is 80--128 for every $k_0$. One run per cell.

\section*{Reviewer 1}

\rev{1.1 Contribution not distinguished from existing frameworks; (j,k)-swap
claim inconsistent.}
The claim that larger release sets give moves ``no swap neighbourhood expresses''
is withdrawn everywhere (it was already withdrawn in Section 3 of the submitted
version and is now absent from the introduction, the abstract and the README). A
region move releasing $j$ vertices and writing back $k$ is a $(j,k)$-swap, and the
text says so. Following the reviewer's observation, Proposition 7 shows
$R(F)=V\setminus N[S\setminus F]$: the region subproblem is the residual MIS
problem with $S\setminus F$ fixed. Table 1 compares release-set selection,
candidate restrictions, objectives, reduction placement and sub-solvers across
Andrade et al., Hashimoto et al., Borowitz et al.\ and this work. Computational
evidence for the additions that can be isolated in one implementation is the
exact/heuristic and reductions/no-reductions ablation (Table 20).

\rev{1.2(a) Reconcile counts and reference sets.}
The abstract's 3/13/11 was against the best of the six \emph{other} solvers, which
is correct and unchanged (16 attain the best of all seven). Table 7 gave 17 because
the table had been built from five solvers; with all seven it is 16. The
ReduMIS-only comparison has 26 pairs (ReduMIS returned nothing on
\texttt{soc-pokec}) and is 3/14/9. Section 5.4 now defines the three reference
sets explicitly. On relative gaps, as the text had always claimed, the Wilcoxon
$p$ against ReduMIS is 0.127 (Holm 0.38), not 0.195.

\rev{1.2(b) Benchmark composition.}
Added: six independent generator seeds at $2^{16}$ and $2^{18}$ and the public
DIMACS10 \texttt{delaunay\_n16/18/20} (Table 21, Section 5.13). \textbf{Baselines
were not run on them}, so the experiment cannot test whether the margin
generalises. It shows that Cascade's value varies by 0.25\% ($2^{16}$) and 0.10\%
($2^{18}$) between graphs against at most 0.05\% between solver seeds, several
times the margin claimed on the seed-1 graph (0.03\%, 0.09\%). The abstract and the
conclusion now restrict the claim to one sample. No additional sparse
large-kernel family was added; I know of no public one and say so.

\rev{1.3 Residual kernel size and direct exact solving.}
Table 14 (Section 5.5): capped kernel, fixed-point kernel, and a 120\,s exact
solve of the whole input, for 19 of the 27 instances (BHOSLIB instances could not
be rebuilt here). Findings: seven instances reduce to $\le126$ vertices and are
proved in $\le1.3$\,s, so region search is irrelevant there; twelve keep
between 0.3\% and 85\% and none is proved by the exact engine in 120\,s. Spearman
correlation of retained fraction with margin is 0.14 ($p=0.56$); the premise of
the paper (large kernel) holds on 12 of 19 instances but does not predict wins:
\texttt{soc-pokec} has a 73\% kernel and is a loss. The submitted text's
``\texttt{com-lj} leaves 9780'' came from another tool; our fixed-point kernel
is 1781, and under the 18\,s cap it ranged from 3544 to 899\,607 across three
runs. The paper now says that the capped residual depends on speed.

\rev{1.4 Ablation does not isolate exact repair and regional reductions.}
New switches \texttt{--lns-heuristic} (reductions then one greedy pass) and
\texttt{--lns-no-reduce}, same implementation, region selection and budget,
three seeds, six instances, with process data (region size before/after the
root reduction, move rate). On \texttt{del18}: $-143$ (heuristic), $-210$ (no
regional reductions), $-1018$ (no region moves). Reductions let regions be
$89$--$203$ vertices that collapse to $5$--$17$; without them $k$ settles at
16--32. A caveat is stated: the heuristic variant's $k$ rule has no proof signal.

\rev{1.5 Recent task-matched solver (InfVC).}
\textbf{Not done.} It is cited and discussed, and Section 5.15 states why the
comparison is absent and that published numbers cannot substitute.

\rev{Minor 1. Provenance and preprocessing.}
\texttt{soc-pokec} and \texttt{com-lj} are relabelled SNAP in the table and the
generator script. The isolated-vertex statement is corrected: removing $q$
isolated vertices lowers $\alpha$ by $q$. SHA-256 checksums of every downloaded
and generated file are in \texttt{results/instance\_checksums.sha256}.

\rev{Minor 2. Cross-platform conclusion.}
Restricted to what was tested, and extended: all 519 solver pairs with a value
on both platforms were compared; 498 have the same sign, 21 change between equal
and strictly ordered, and none reverses.

\section*{Reviewer 2}

\rev{Main idea hard to follow.}
The introduction now opens the idea with a plain example (the $(1,2)$-swap, then
the same question for a set of released vertices) before any notation, and refers
to the worked figure.

\rev{Relations to existing work.}
Katzmann and Komusiewicz (AAAI 2017) is added, with Hashimoto et al.\ and
Borowitz et al.

\rev{Comparison to state-of-the-art solvers; clearer conclusion.}
Abstract and Conclusions state the result as it is: strictly ahead on three
Delaunay graphs from one seed, 13 ties, 11 losses, no significant difference from
ReduMIS, NuMVC or FastVC, third of seven by mean gap. See ``What remains open''
for what is still missing.

\section*{Reviewer 4}

The reviewer is right about the central defect: methods missing from tables that
the text refers to. The cause was \texttt{make\_tables.py} (a five-solver list and
the wrong default input file). It is fixed, and every table is regenerated from
\texttt{results/cor\_main60.csv}. I then checked the text against the data rather
than only the tables. Additional errors found and corrected:
\begin{itemize}
\item the seed-study significance claims (FastVC $p=0.016$, NuMVC $p=0.023$) came
  from raw sizes; on relative gaps with Holm correction only OnlineMIS is
  significant;
\item ``significantly ahead of five of six baselines'' $\to$ three of six;
\item ``mean gap 0.414\%'' $\to$ 0.785\% (third of seven); FastVC is lowest, 0.340\%;
\item the ablation text now states that two ablated variants beat the shipped
  configuration on \texttt{del16}/\texttt{del18}, within seed noise;
\item the 543\,$\mu$s per move figure now has a record (the new runs reproduce it:
  $560\,\mu$s on \texttt{del18});
\item the queue-ordering timings had no record and were removed;
\item a table of the whence-ReduMIS decomposition listed 27 instances where 26
  were paired.
\end{itemize}
The compile has no undefined references or citations. Repository: the
\texttt{reviews/} directory and committed build outputs were removed.

\section*{What remains open}
\begin{enumerate}
\item \textbf{InfVC and any 2024--2026 solver}: not run. Required by Reviewers 1 and 3.
\item \textbf{Baselines on the independent Delaunay graphs}: not run, so the claim
  that Cascade leads on Delaunay graphs rests on one sample. Running ReduMIS and
  FastVC on Table 21's graphs is the single experiment that decides whether the
  positive result exists.
\item \textbf{BHOSLIB kernel statistics}: not measured (instances not available
  here).
\item \textbf{Cascade vs.\ the Borowitz et al.\ implementation}: not run.
\item The new experiments ran on a different machine (4-core Xeon 2.8\,GHz,
  three concurrent runs) from the main tables; the paper says so.
\item Author list of InfVC in \texttt{refs.bib} is incomplete (first author only
  verified) and the two Gro{\ss}mann et al.\ items are arXiv entries; please
  check the bibliographic details before submission.
\item EJOR's exact limits (abstract length, highlights, structure) were not
  retrievable; the abstract is 147 words and the five highlights are
  $\le 77$ characters.
\end{enumerate}
\end{document}
